\documentclass[11pt]{article}
\pdfinfoomitdate=1
\pdftrailerid{}
\pdfsuppressptexinfo=15
\usepackage[a4paper,margin=1in]{geometry}
\usepackage{amsmath,amssymb,amsthm,booktabs,enumitem,microtype}
\usepackage[numbers,sort&compress]{natbib}
\usepackage{cleveref}

\newtheorem{theorem}{Theorem}[section]
\newtheorem{corollary}[theorem]{Corollary}
\newtheorem{proposition}[theorem]{Proposition}
\theoremstyle{remark}
\newtheorem{remark}[theorem]{Remark}
\newcommand{\Tr}{\operatorname{Tr}}
\newcommand{\cA}{\mathcal A}
\newcommand{\cB}{\mathcal B}
\newcommand{\cE}{\mathcal E}
\newcommand{\cP}{\mathcal P}
\newcommand{\cR}{\mathcal R}
\newcommand{\cS}{\mathcal S}

\title{A First-Lower-Sideband Orbit Atom and a\\
Mandatory Compensation Obstruction}
\author{Bin Wang}
\date{August 2026}

\begin{document}
\maketitle

\begin{abstract}
The one-alias off-background of RH-330 contains every coefficient below the
critical order $2k$ except the critical coefficient itself.  We prove that
the first lower even sideband $n_-=2k-2$ already contains an alias-scale
physical boundary-orbit atom.  The primitive period-$2(k-1)$ boundary orbit
has a folded subset of $2k-3$ marked points in the frozen RH-334 far set.
Its noisy localized trace is exactly zero, while its deterministic mass
$D_{k-1}^{\rm orb}$ satisfies
\[
 D_{k-1}^{\rm orb}/H_{k-1}\longrightarrow+\infty,
 \qquad H_{k-1}=(k-1)R^{-2(k-1)}.
\]
Writing the actual typed sideband coefficient as
$q_-= -D_{k-1}^{\rm orb}+C_-$, vanishing of the off-alias weighted prefix
necessarily forces
$C_-=D_{k-1}^{\rm orb}+o(H_{k-1})$, at relative precision
$o((\beta R)^{-2(k-1)})$.  Thus any proof that takes separate absolute
values after the orbit/complement split fails at this mandatory sideband.
This is not a positive lower bound for the fully signed coefficient: the
complement may cancel the atom.  The repository provides no moving-order
estimate for that complement, so both off-alias vanishing and off-alias
nonvanishing remain not testable.  The physical counterloop radial sideband
is retained without a sign claim because the printed value $C_M>1$ lacks an
interval certificate.  No determinant or Riemann-hypothesis conclusion is
obtained.
\end{abstract}

\section{Typed off-alias coefficient}

Work on the physical natural clock
\begin{equation}\label{eq:clock}
 k=k_\sigma\in\mathbb N,\qquad k\ge2,\qquad
 k=\frac{\log(1/\sigma)}{2\log\lambda}+O(1),
 \qquad \sigma\to0,
 \qquad H_k=kR^{-2k},
 \qquad R=\frac75.
\end{equation}
Choose any integer cut
\begin{equation}\label{eq:cut}
 2k<h_\sigma\le4k.
\end{equation}
For the actual Hardy full-trace constituent $q_{\sigma,k,n}$, RH-330 defines
the nonnegative off-alias prefix
\begin{equation}\label{eq:E-off}
 \cE_{\sigma,\rm off}^{(h)}(R)
 =\sum_{\substack{2\le n<h_\sigma\\ n\ne2k}}
 \frac{|q_{\sigma,k,n}|R^n}{n}
\end{equation}
\citep{WangTransfer2026}.

We retain the corrected RH-334 folded backward-observable operator
$K_\sigma$, the frozen basepoint partition
$[0,1]=J^-\sqcup J^+\sqcup F$, and the localized raw slots
$\cB_{\sigma,k,n},\cS_{\sigma,k,n},\cR_{\sigma,k,n}$
\citep{WangObservation2026}.  For every $n\ge2$, define
\begin{align}
 q_{\sigma,k,n}
 &=c^H_{\sigma,n}-s_{k,n}-a_n^{\rm num},
 \label{eq:q-definition}\\
 \cP_{\sigma,n}
 &=r_H^{-n}\{(-1)^n-\lambda_-(\sigma)^n\},
 \qquad r_H=\frac{17}{20},
 \label{eq:parity}\\
 \cA_{k,n}&=s_{k,n}-p_n^{\rm pole}.
 \label{eq:counterloop-defect}
\end{align}

\begin{proposition}[All-order typed five-slot identity]
\label{prop:five-slot}
For every $n\ge2$,
\begin{equation}\label{eq:five-slot}
 \boxed{
 q_{\sigma,k,n}
 =\cB_{\sigma,k,n}+\cS_{\sigma,k,n}+\cR_{\sigma,k,n}
  +\cP_{\sigma,n}-\cA_{k,n}.}
\end{equation}
The coefficient type is the Hardy full-trace constituent, not automatically
the modulus-complement coefficient.
\end{proposition}

\begin{proof}
The deterministic numerator anchor is
$a_n^{\rm num}=c^H_n-p_n^{\rm pole}$.  Hence
\begin{equation*}
 q_{\sigma,k,n}
 =(c^H_{\sigma,n}-c^H_n)-(s_{k,n}-p_n^{\rm pole}).
\end{equation*}
The exact bulk split gives
\begin{equation*}
 c^H_{\sigma,n}-c^H_n
 =r_H^{-n}\{\Tr K_\sigma^n-P_n\}+\cP_{\sigma,n}.
\end{equation*}
RH-334 partitions the raw trace into $\cB+\cS+\cR$ for every $n\ge2$.
Substitution proves \eqref{eq:five-slot}.  Identifying this constituent with
the modulus complement would additionally require the noisy-head defect to
vanish, which is not assumed here.
\end{proof}

Set
\begin{equation}\label{eq:n-minus}
 n_-=2k-2=2m,
 \qquad m=k-1,
 \qquad H_m=mR^{-2m}.
\end{equation}
For every admissible cut \eqref{eq:cut}, $2\le n_-<h_\sigma$ and
$n_-\ne2k$.  Therefore the corresponding summand is always present in
\eqref{eq:E-off}.

At this sideband the exact RH-326 counterloop identity is
\begin{equation}\label{eq:radial-sideband}
 \cA_{k,n_-}=2(\beta^{n_-}-\beta_k^{n_-}),
 \qquad
 \beta_k=\frac{|M_k|^{-1/(2k)}}{r_H},
 \qquad
 \beta=\frac1{r_H\sqrt\lambda}
\end{equation}
\citep{WangFirstAlias2026}.  RH-17 proves only $C_M>0$ analytically in
$|M_k|=C_M\lambda^k(1+o(1))$.  Its printed value
$1.9463429052\ldots$ is not an interval certificate, so no eventual sign of
\eqref{eq:radial-sideband} is used.  Here ``alias-scale'' refers only to the
size of the orbit atom.  The lower sideband is not an alias impulse of the
current period-$2k$ counterloop; \eqref{eq:radial-sideband} is its radial
counterloop contribution.

\section{A physical orbit atom at $n_-=2k-2$}

Let $p_{2m}$ be the primitive boundary point of physical period $2m$ and
write $x_j=f^j(p_{2m})$.  RH-17 supplies the ordered component chain, fixed
gaps from the fold cusp $b$, and the multiplier law
\begin{equation}\label{eq:M-m}
 M_m=(f^{2m})'(p_{2m})=-C_M\lambda^m\{1+o(1)\}
\end{equation}
\citep{WangBoundaryMonodromy2026}.  Define
\begin{equation}\label{eq:Omega-m}
 \Omega_m=\{|x_j|:0\le j<2m,\ j\ne2m-2\}.
\end{equation}

\begin{theorem}[Mandatory lower-sideband far atom]
\label{thm:sideband-atom}
For every fixed RH-334 window parameter $A>0$ and all sufficiently large
$k$,
\begin{equation}\label{eq:Omega-sideband-F}
 \Omega_m\subset F_{\sigma,k,A},
 \qquad |\Omega_m|=2m-1=2k-3.
\end{equation}
Define
\begin{equation}\label{eq:D-m}
 D_m^{\rm orb}
 =r_H^{-2m}\frac{2m-1}{1+|M_m|}>0.
\end{equation}
Then the localized sideband far subledger is exactly
\begin{equation}\label{eq:R-orb-minus}
 \cR_{k,-}^{\rm orb}
 :=r_H^{-2m}\{L_{\sigma,2m}(\Omega_m)
              -P_{2m}^{\rm abs}(\Omega_m)\}
 =-D_m^{\rm orb},
\end{equation}
and
\begin{equation}\label{eq:D-Hm}
 \frac{D_m^{\rm orb}}{H_m}
 =\frac{2}{C_M}(\beta R)^{2m}\{1+o(1)\}
 \longrightarrow+\infty.
\end{equation}
\end{theorem}

\begin{proof}
The fixed-gap proof of RH-338 applies to the period-$2m$ orbit even though
the windows are evaluated on the narrower $k$-clock
\citep{WangFarAtom2026}.  The endpoint lies to the right of $b$ by a fixed
gap; all retained internal even points are at most $h(b)<b$; and all odd
folded points are at most $r<b$.  Since $A\sqrt\sigma\to0$, deleting only
$h(p_{2m})$ leaves $2m-1$ distinct points in $F_{\sigma,k,A}$.

The finite set $\Omega_m$ has zero multiplication operator on $L^2$, so its
noisy localized trace is zero.  Folding preserves the multiplier and marked
multiplicity.  Since $M_m<0$ for large $m$, each deterministic weight is
$1/(1+|M_m|)$, proving \eqref{eq:R-orb-minus}.  Finally
\eqref{eq:M-m} and $\beta^{2m}=r_H^{-2m}\lambda^{-m}$ give
\begin{equation*}
 D_m^{\rm orb}=\frac{2m}{C_M}\beta^{2m}\{1+o(1)\}.
\end{equation*}
Division by $H_m$ gives \eqref{eq:D-Hm}; RH-336 proves $\beta R>1$ exactly
\citep{WangProjectorMass2026}.
\end{proof}

As in RH-338, $\Omega_m$ is only an analytic subpartition of the already
frozen far set.  It is not a new canonical physical observation window.

\section{Off-alias compensation obstruction}

Collect every other signed contribution at $n_-$ into the exact complement
\begin{equation}\label{eq:C-minus}
 C_k^-:=q_{\sigma,k,n_-}+D_m^{\rm orb}.
\end{equation}
Thus
\begin{equation}\label{eq:q-minus-split}
 q_{\sigma,k,n_-}=-D_m^{\rm orb}+C_k^-.
\end{equation}
$C_k^-$ includes the rest of the raw far slot, the boundary and sibling
slots, the parity packet, and the radial counterloop term
\eqref{eq:radial-sideband}, with their actual signs retained.

\begin{theorem}[Necessary compensation for off-alias vanishing]
\label{thm:necessary-compensation}
For every admissible cut \eqref{eq:cut},
\begin{equation}\label{eq:E-lower-coeff}
 \cE_{\sigma,\rm off}^{(h)}(R)
 \ge\frac{|q_{\sigma,k,n_-}|R^{2m}}{2m}
 =\frac{|q_{\sigma,k,n_-}|}{2H_m}.
\end{equation}
Consequently,
\begin{equation}\label{eq:E-implies-compensation}
 \cE_{\sigma,\rm off}^{(h)}(R)\longrightarrow0
 \quad\Longrightarrow\quad
 C_k^-=D_m^{\rm orb}+o(H_m).
\end{equation}
The required relative precision is
\begin{equation}\label{eq:relative-sideband}
 o\!\left(\frac{H_m}{D_m^{\rm orb}}\right)
 =o\!\left((\beta R)^{-2m}\right).
\end{equation}
Moreover, an atom/complement proof that takes separate absolute values
before their signed sum cannot close, because the orbit atom alone contributes
\begin{equation}\label{eq:absolute-sideband}
 \frac{D_m^{\rm orb}R^{2m}}{2m}
 =\frac{D_m^{\rm orb}}{2H_m}
 \longrightarrow+\infty.
\end{equation}
\end{theorem}

\begin{proof}
The $n_-$ summand is present in \eqref{eq:E-off}, proving
\eqref{eq:E-lower-coeff}.  If the nonnegative prefix tends to zero, that
summand tends to zero and hence
$q_{\sigma,k,n_-}=o(H_m)$.  Substitute \eqref{eq:q-minus-split} to obtain
\eqref{eq:E-implies-compensation}.  The reciprocal of \eqref{eq:D-Hm}
gives \eqref{eq:relative-sideband}; \eqref{eq:absolute-sideband} is the same
identity without the signed complement.
\end{proof}

This theorem is one-way.  It does not lower-bound the fully signed
$|q_{\sigma,k,n_-}|$, because $C_k^-$ may compensate the orbit atom.  RH-19
leaves the requisite bulk moving-order control open
\citep{WangComplement2026}.  Therefore neither
$\cE_{\sigma,\rm off}^{(h)}(R)\to0$ nor its failure follows from the current
repository.  The exact verdict is a scoped obstruction to the
separate-absolute sideband route and
\texttt{NOT\_TESTABLE} for the aggregate off-alias background.

\section{Reproduction and claim boundary}

The artifact evaluates the physical algebraic constants and period-$2m$
boundary orbits at high Decimal precision for
$k=3,5,9,17,33$ and $A=1/4$.  It reproduces the mandatory sideband orders,
far counts $3,7,15,31,63$, orbit closure, the identity
$R^{2m}/(2m)=1/(2H_m)$, and the super-target orbit mass.  These finite rows
are diagnostic formula checks, not interval certificates or moving-order
evidence.

The paper proves only the actual typed lower-sideband atom and the necessary
signed compensation condition.  It proves neither off-alias vanishing nor a
positive lower bound for the fully signed sideband coefficient.  It does not
close the critical packet, noisy head, counterloop synchronization,
determinant gluing, full-trace replacement, or full-trace divergence.  No
Hilbert--Polya operator, Riemann-zero identification, von Mangoldt trace,
completed-zeta divisor equality, or Riemann-hypothesis proof is obtained.
Gates A--E remain false/open.

\bibliographystyle{plainnat}
\bibliography{references}

\end{document}
